
\documentclass{article}
\usepackage{colortbl}
\usepackage{makecell}
\usepackage{multirow}
\usepackage{supertabular}

\begin{document}

\newcounter{utterance}

\centering \large Interaction Transcript for game `cladder', experiment `full\_v1.5\_default', episode 9587 with Qwen3.5{-}9B.
\vspace{24pt}

{ \footnotesize  \setcounter{utterance}{1}
\setlength{\tabcolsep}{0pt}
\begin{supertabular}{c@{$\;$}|p{.15\linewidth}@{}p{.15\linewidth}p{.15\linewidth}p{.15\linewidth}p{.15\linewidth}p{.15\linewidth}}
    \# & \multicolumn{2}{c}{Player} && \multicolumn{2}{c}{Game Master} \\
    \hline

    \theutterance \stepcounter{utterance}  
    & & & \multicolumn{4}{p{0.6\linewidth}}{
        \cellcolor[rgb]{0.9,0.9,0.9}{
            \makecell[{{p{\linewidth}}}]{
                \texttt{\tiny{[P1$\langle$GM]}}
                \texttt{You are an expert in causal inference. The following question is not a typical commonsense query, but rather a meticulously designed question created by a professor specializing in causal inference, intended to assess the students' mastery of the course content.} \\
\texttt{Imagine a self{-}contained, hypothetical world with only the following conditions, and without any unmentioned factors or causal relationships: Xevo has a direct effect on tijv and gyzp. Tijw has a direct effect on tijv and gyzp. Tijv has a direct effect on gyzp. Tijw is unobserved. We know that xevo or tijw causes tijv. xevo or tijw or tijv causes gyzp. We observed an individual is tijw. Would an individual is not gyzp if not xevo instead of xevo?} \\
\texttt{Start your answer with 'Yes' or 'No', followed by additional reasoning or evidence to support your explanation.} \\
            }
        }
    }
     \\ \\

    \theutterance \stepcounter{utterance}  
    & \multicolumn{4}{p{0.6\linewidth}}{
        \cellcolor[rgb]{0.561,0.737,0.561}{
            \makecell[{{p{\linewidth}}}]{
                \texttt{\tiny{[P1$\rangle$GM]}}
                \texttt{No} \\
\\ 
\texttt{**Reasoning:**} \\
\\ 
\texttt{To determine the answer, we must analyze the causal structure and the specific counterfactual query within the constraints of the hypothetical world.} \\
\\ 
\texttt{1.  **Causal Structure Analysis:**} \\
\texttt{    *   **Parents of `tijv`:** The prompt states "Xevo has a direct effect on `tijv`" and "Tijw has a direct effect on `tijv`." It further clarifies that "xevo or tijw causes tijv." Thus, \$Parents(tijv) = \textbackslash{}\{xevo, tijw\textbackslash{}\}\$.} \\
\texttt{    *   **Parents of `gyzp`:** The prompt states "Xevo has a direct effect on `gyzp`," "Tijw has a direct effect on `gyzp`," and "Tijv has a direct effect on `gyzp`." Thus, \$Parents(gyzp) = \textbackslash{}\{xevo, tijw, tijv\textbackslash{}\}\$.} \\
\texttt{    *   **Observed State:** We observe an individual where \$tijw = 1\$ (is `tijw`). The value of `xevo` is not explicitly fixed in the observation, but the query asks about the outcome under a specific intervention.} \\
\\ 
\texttt{2.  **The Counterfactual Query:**} \\
\texttt{    The question asks: "Would an individual {[}be{]} not `gyzp` if not `xevo` instead of `xevo`?"} \\
\texttt{    Formally,} \\
            }
        }
    }
    & & \\ \\

    \theutterance \stepcounter{utterance}  
    & & & \multicolumn{2}{p{0.3\linewidth}}{
        \cellcolor[rgb]{0.9,0.9,0.9}{
            \makecell[{{p{\linewidth}}}]{
                \texttt{\tiny{[GM$|$GM]}}
                \texttt{No} \\
            }
        }
    }
    & & \\ \\

    \theutterance \stepcounter{utterance}  
    & & & \multicolumn{2}{p{0.3\linewidth}}{
        \cellcolor[rgb]{0.9,0.9,0.9}{
            \makecell[{{p{\linewidth}}}]{
                \texttt{\tiny{[GM$|$GM]}}
                \texttt{no} \\
            }
        }
    }
    & & \\ \\

    \theutterance \stepcounter{utterance}  
    & & & \multicolumn{2}{p{0.3\linewidth}}{
        \cellcolor[rgb]{0.9,0.9,0.9}{
            \makecell[{{p{\linewidth}}}]{
                \texttt{\tiny{[GM$|$GM]}}
                \texttt{game\_result = WIN} \\
            }
        }
    }
    & & \\ \\

\end{supertabular}
}

\end{document}
